\documentclass[11pt,a4paper]{article}
\usepackage[utf8]{inputenc}
\usepackage{graphicx}
\usepackage{multirow}
\usepackage{enumitem}
\usepackage{amsmath,amssymb}
\usepackage{array}
\usepackage[left=10mm,right=20mm,top=20mm,bottom=10mm]{geometry}

\newcommand{\R}{\mathbb{R}}
\newcommand{\Span}{\operatorname{span}}
\graphicspath{{C:/Users/teoso/Documents/Template Math Depart/}}
\setlist[enumerate,1]{label=\arabic*.,leftmargin=*,itemsep=9pt,parsep=0pt,topsep=6pt}
\setlist[enumerate,2]{label=(\alph*),leftmargin=*,itemsep=6pt,parsep=0pt,topsep=5pt}

\begin{document}
\pagenumbering{gobble}
\setlength{\abovedisplayskip}{7pt}
\setlength{\belowdisplayskip}{7pt}
\setlength{\abovedisplayshortskip}{4pt}
\setlength{\belowdisplayshortskip}{4pt}

\noindent
\begin{minipage}{0.15\textwidth}
  \includegraphics[width=1.8cm]{ITS.png}
\end{minipage}%
\begin{minipage}{0.7\textwidth}
  \centering
  \textbf{EVALUASI TENGAH SEMESTER GANJIL 2025/2026}\\
  \textbf{DEPARTEMEN MATEMATIKA -- FSAD ITS}
\end{minipage}%
\begin{minipage}{0.15\textwidth}
  \centering
  \includegraphics[width=1.8cm]{M.png}
\end{minipage}
\begin{center}
  \begin{tabular}{@{}l c l@{}}
    Mata kuliah / Kode & : & Aljabar Linier                         \\
    Hari / Tanggal     & : & Kamis, 16 Oktober 2025                 \\
    Sifat / waktu      & : & Tertutup / 100 menit                   \\
    Dosen              & : & Prof. Subiono, M.S.                    \\
                       &   & Dr. Dian Winda Setyawati, S.Si., M.Si. \\
                       &   & Soleha, S.Si., M.Si.                   \\
                       &   & Dr. Dieky Adzkiya, S.Si., M.Si.
  \end{tabular}
\end{center}

\noindent\textbf{HARAP DIPERHATIKAN !!!}\\
\textit{Segala jenis pelanggaran (mencontek, kerjasama, dsb) yang dilakukan saat ETS/EAS
  akan dikenakan sanksi pembatalan mata kuliah pada semester yang sedang berjalan.}

\vspace{0.3cm}
\noindent ETS mengukur kemampuan
\begin{center}
  \begin{tabular}{|>{\centering\arraybackslash}p{1.6cm}|p{\dimexpr\textwidth-5.55cm-12\tabcolsep-7\arrayrulewidth\relax}|>{\centering\arraybackslash}p{2cm}|>{\centering\arraybackslash}p{0.65cm}|>{\centering\arraybackslash}p{0.65cm}|>{\centering\arraybackslash}p{0.65cm}|}
    \hline
    \multicolumn{2}{|c|}{CPMK}
                           & \multirow{2}{*}{\shortstack{Bobot\\CPMK (\%)}}
                           & \multicolumn{3}{c|}{CPL}                                                                                                                              \\ \cline{4-6}
    \multicolumn{2}{|c|}{} &                                                                                                          & 4            & 5            & 9            \\ \hline
    CPMK-1                 & Mahasiswa menguasai dan mampu menjelaskan secara tertulis mengenai konsep ruang vektor secara umum
                           & 5                                                                                                        & $\checkmark$ & $\checkmark$ & $\checkmark$ \\ \hline
    CPMK-2                 & Mahasiswa menguasai dan mampu menjelaskan secara tertulis mengenai transformasi linier pada ruang vektor
                           & 10                                                                                                       & $\checkmark$ & $\checkmark$ & $\checkmark$ \\ \hline
    CPMK-3                 & Mahasiswa menguasai dan mampu menerapkan proses diagonalisasi atau dekomposisi suatu matriks
                           & 5                                                                                                        & $\checkmark$ & $\checkmark$ & $\checkmark$ \\ \hline
  \end{tabular}
\end{center}

\begin{center}
  \begin{tabular}{|c|c|c|c|c|}
    \hline
    \multirow{2}{*}{Soal}                  & \multirow{2}{*}{Bobot Soal (\%)} & \multicolumn{3}{c|}{}                   \\ \cline{3-5}
                                           &                                  & CPMK-1                & CPMK-2 & CPMK-3 \\ \hline
    1                                      & 25                               & 5                     &        &        \\ \hline
    2                                      & 25                               &                       & 5      &        \\ \hline
    3                                      & 25                               &                       & 5      &        \\ \hline
    4                                      & 25                               &                       &        & 5      \\ \hline
    \multicolumn{2}{|c|}{Total Bobot CPMK} & 5                                & 10                    & 5               \\ \hline
  \end{tabular}
\end{center}

\begin{center}
  Soal
\end{center}
\begin{enumerate}[series=soal]
  \item Diberikan ruang vektor $M_{2\times2}(\R)$ atas lapangan $\R$,
        dan $S\subseteq M_{2\times2}(\R)$ dengan
        \[
          S=\left\{
          \begin{bmatrix}2&-1\\1&3\end{bmatrix},
          \begin{bmatrix}0&4\\-1&2\end{bmatrix},
          \begin{bmatrix}4&2\\1&3\end{bmatrix}
          \right\}.
        \]
        \begin{enumerate}
          \item Apakah himpunan $S$ bebas linier? Jelaskan jawaban Saudara.
          \item Apakah $S$ basis untuk $M_{2\times2}(\R)$? Jelaskan jawaban Saudara.
        \end{enumerate}

  \item Diberikan $T:P_3(\R)\to\R^4$ didefinisikan oleh
        \[
          T(a+bx+cx^2+dx^3)=\begin{bmatrix}a+b\\a-b\\c\\c+a\end{bmatrix}.
        \]
        \begin{enumerate}
          \item Tentukan $T(1+2x+3x^2+4x^3)$.
          \item Tunjukkan bahwa $T$ transformasi linear.
          \item Tentukan basis dari $N(T)$ dan dimensinya.
          \item Tentukan basis dari $R(T)$ dan dimensinya.
        \end{enumerate}
\end{enumerate}

\vfill
\begin{center}
  \textit{Selamat Mengerjakan}
\end{center}
\newpage

\begin{enumerate}[resume=soal]
  \item Diberikan $V,W$ ruang vektor atas lapangan $F$ dan $T$ transformasi linear $T:V\to W$.
        Basis untuk $V$ adalah $B=\{v_1,v_2\}$.
        Basis untuk $W$ adalah $C=\{w_1,w_2,w_3\}$.
        Misalkan $T(v_1)=2w_1+3w_2+4w_3$ dan $T(v_2)=w_1+2w_3$.
        \begin{enumerate}
          \item Tentukan matriks representasi $[T]_B^C$.
          \item Tentukan $[T(v)]_C$ jika $[v]_B=\begin{bmatrix}1\\-2\end{bmatrix}$, $v\in V$.
        \end{enumerate}

  \item Diberikan matriks $A_{3\times3}$
        \[
          A=\begin{bmatrix}4&0&1\\2&3&2\\1&0&4\end{bmatrix}.
        \]
        \begin{enumerate}
          \item Buktikan bahwa persamaan karakteristik dari $A$ adalah $(3-\lambda)^2(5-\lambda)=0$.
          \item Dapatkan ruang-ruang eigen matriks $A$.
          \item Dapatkan matriks nonsingular $P$ sedemikian hingga $D=P^{-1}AP$ dengan $D$ matriks diagonal.
        \end{enumerate}
\end{enumerate}

\vfill
\begin{center}
  \textit{Selamat Mengerjakan}
\end{center}
\newpage

\begin{center}
  Solusi
\end{center}
\begin{enumerate}[itemsep=12pt,parsep=5pt]
  \item Namai ketiga matriks pada $S$ sebagai $A_1,A_2,A_3$, sesuai urutannya.
        \begin{enumerate}
          \item Untuk memeriksa kebebasan linier, misalkan
                \[
                  \alpha A_1+\beta A_2+\gamma A_3=0.
                \]
                Samakan setiap entri dengan entri matriks nol.
                Diperoleh sistem
                \[
                  \begin{cases}
                    2\alpha+4\gamma=0,        \\
                    -\alpha+4\beta+2\gamma=0, \\
                    \alpha-\beta+\gamma=0,    \\
                    3\alpha+2\beta+3\gamma=0.
                  \end{cases}
                \]
                Persamaan pertama memberikan $\alpha=-2\gamma$.
                Substitusi ke persamaan ketiga memberikan
                \[
                  -2\gamma-\beta+\gamma=0
                  \quad\Longrightarrow\quad\beta=-\gamma.
                \]
                Masukkan keduanya ke persamaan keempat:
                \[
                  3(-2\gamma)+2(-\gamma)+3\gamma=0
                  \quad\Longrightarrow\quad-5\gamma=0.
                \]
                Jadi, $\gamma=0$, lalu $\alpha=\beta=0$.
                Hanya kombinasi linear trivial yang menghasilkan matriks nol.
                Maka, $S$ bebas linier.

          \item Ruang $M_{2\times2}(\R)$ berdimensi $4$.
                Setiap matriksnya memiliki empat entri yang dapat dipilih bebas.
                Jadi, setiap basis ruang ini harus memiliki empat anggota.

                Himpunan $S$ hanya memiliki tiga matriks.
                Karena itu, $S$ tidak merentang seluruh $M_{2\times2}(\R)$.
                Maka, $S$ bukan basis untuk $M_{2\times2}(\R)$.
        \end{enumerate}

  \item Untuk $p(x)=a+bx+cx^2+dx^3$, nilai $T(p)$ ditentukan oleh koefisien $a,b,c$.
        \begin{enumerate}
          \item Pada $p(x)=1+2x+3x^2+4x^3$, diperoleh $a=1$, $b=2$, $c=3$, dan $d=4$.
                Dengan demikian,
                \[
                  T(1+2x+3x^2+4x^3)
                  =\begin{bmatrix}1+2\\1-2\\3\\3+1\end{bmatrix}
                  =\begin{bmatrix}3\\-1\\3\\4\end{bmatrix}.
                \]

          \item Ambil $p(x)=a+bx+cx^2+dx^3$ dan $q(x)=e+fx+gx^2+hx^3$.
                Untuk skalar $\alpha,\beta\in\R$, berlaku
                \begin{align*}
                  T(\alpha p+\beta q)
                   & =\begin{bmatrix}
                        \alpha(a+b)+\beta(e+f) \\
                        \alpha(a-b)+\beta(e-f) \\
                        \alpha c+\beta g       \\
                        \alpha(c+a)+\beta(g+e)
                      \end{bmatrix}                           \\
                   & =\alpha\begin{bmatrix}a+b\\a-b\\c\\c+a\end{bmatrix}
                  +\beta\begin{bmatrix}e+f\\e-f\\g\\g+e\end{bmatrix}     \\
                   & =\alpha T(p)+\beta T(q).
                \end{align*}
                Jadi, $T$ memenuhi sifat linear.

          \item Ruang nol $N(T)$ berisi semua polinomial yang dipetakan ke vektor nol.
                Syarat $T(p)=0$ memberikan
                \[
                  a+b=0,\qquad a-b=0,\qquad c=0,\qquad c+a=0.
                \]
                Jumlahkan dua persamaan pertama.
                Diperoleh $2a=0$, sehingga $a=0$.
                Maka, $b=0$ dan $c=0$.
                Koefisien $d$ tetap bebas karena tidak muncul dalam $T(p)$.
                Jadi, setiap anggota $N(T)$ berbentuk $dx^3$ dengan $d\in\R$.
                Polinomial $x^3$ bukan polinomial nol, sehingga $\{x^3\}$ bebas linier.
                Maka,
                \[
                  N(T)=\Span\{x^3\},\qquad B_{N(T)}=\{x^3\},\qquad\dim N(T)=1.
                \]

          \item Ruang hasil $R(T)$ berisi semua keluaran $T$.
                Tuliskan keluaran umumnya sebagai kombinasi linear:
                \[
                  T(p)=\begin{bmatrix}a+b\\a-b\\c\\c+a\end{bmatrix}
                  =a\begin{bmatrix}1\\1\\0\\1\end{bmatrix}
                  +b\begin{bmatrix}1\\-1\\0\\0\end{bmatrix}
                  +c\begin{bmatrix}0\\0\\1\\1\end{bmatrix}.
                \]
                Semua nilai $a,b,c\in\R$ dapat dipilih bebas.
                Maka, ketiga vektor tersebut merentang $R(T)$.

                Untuk memeriksa kebebasan liniernya, misalkan
                \[
                  r\begin{bmatrix}1\\1\\0\\1\end{bmatrix}
                  +s\begin{bmatrix}1\\-1\\0\\0\end{bmatrix}
                  +t\begin{bmatrix}0\\0\\1\\1\end{bmatrix}=0.
                \]
                Dua komponen pertama memberikan $r+s=0$ dan $r-s=0$.
                Jadi, $r=s=0$.
                Komponen ketiga memberikan $t=0$.
                Maka, ketiganya bebas linier.
                Dengan demikian,
                \[
                  B_{R(T)}=\left\{
                  \begin{bmatrix}1\\1\\0\\1\end{bmatrix},
                  \begin{bmatrix}1\\-1\\0\\0\end{bmatrix},
                  \begin{bmatrix}0\\0\\1\\1\end{bmatrix}
                  \right\},
                  \qquad\dim R(T)=3.
                \]
                Hasil ini sesuai dengan $\dim N(T)+\dim R(T)=1+3=4=\dim P_3(\R)$.
        \end{enumerate}

  \item Kolom matriks representasi diambil dari koordinat bayangan setiap vektor basis domain.
        \begin{enumerate}
          \item Dari bentuk $T(v_1)$ dan $T(v_2)$, diperoleh
                \[
                  [T(v_1)]_C=\begin{bmatrix}2\\3\\4\end{bmatrix},
                  \qquad
                  [T(v_2)]_C=\begin{bmatrix}1\\0\\2\end{bmatrix}.
                \]
                Susun keduanya sebagai kolom sesuai urutan $v_1,v_2$.
                Matriks representasinya adalah
                \[
                  [T]_B^C
                  =\begin{bmatrix}[T(v_1)]_C&[T(v_2)]_C\end{bmatrix}
                  =\begin{bmatrix}2&1\\3&0\\4&2\end{bmatrix}.
                \]

          \item Gunakan hubungan koordinat $[T(v)]_C=[T]_B^C[v]_B$.
                Maka,
                \[
                  [T(v)]_C
                  =\begin{bmatrix}2&1\\3&0\\4&2\end{bmatrix}
                  \begin{bmatrix}1\\-2\end{bmatrix}
                  =\begin{bmatrix}2-2\\3+0\\4-4\end{bmatrix}
                  =\begin{bmatrix}0\\3\\0\end{bmatrix}.
                \]
                Artinya, $T(v)=3w_2$.
                Ini juga diperoleh dari $v=v_1-2v_2$ dan linearitas $T$.
        \end{enumerate}

  \item
        \begin{enumerate}
          \item Persamaan karakteristik diperoleh dari $\det(A-\lambda I)=0$.
                Ekspansikan determinan pada kolom kedua:
                \begin{align*}
                  \det(A-\lambda I)
                   & =\begin{vmatrix}4-\lambda&0&1\\2&3-\lambda&2\\1&0&4-\lambda\end{vmatrix} \\
                   & =(3-\lambda)\begin{vmatrix}4-\lambda&1\\1&4-\lambda\end{vmatrix}         \\
                   & =(3-\lambda)\bigl((4-\lambda)^2-1\bigr)                                  \\
                   & =(3-\lambda)(3-\lambda)(5-\lambda)                                       \\
                   & =(3-\lambda)^2(5-\lambda).
                \end{align*}
                Jadi, persamaan karakteristiknya adalah $(3-\lambda)^2(5-\lambda)=0$.
                Nilai eigennya adalah $\lambda=3$ dan $\lambda=5$.

          \item Untuk $\lambda=3$, selesaikan $(A-3I)v=0$.
                Tuliskan $v=(x,y,z)^{\mathsf T}$.
                Karena
                \[
                  A-3I=\begin{bmatrix}1&0&1\\2&0&2\\1&0&1\end{bmatrix},
                \]
                semua persamaannya setara dengan $x+z=0$.
                Maka, $z=-x$ dan $y$ bebas.
                Setiap solusinya berbentuk
                \[
                  v=\begin{bmatrix}x\\y\\-x\end{bmatrix}
                  =x\begin{bmatrix}1\\0\\-1\end{bmatrix}
                  +y\begin{bmatrix}0\\1\\0\end{bmatrix}.
                \]
                Jadi,
                \[
                  E_3=\Span\left\{
                  \begin{bmatrix}1\\0\\-1\end{bmatrix},
                  \begin{bmatrix}0\\1\\0\end{bmatrix}
                  \right\},\qquad\dim E_3=2.
                \]

                Untuk $\lambda=5$, selesaikan $(A-5I)v=0$.
                Sistemnya adalah
                \[
                  \begin{cases}
                    -x+z=0,     \\
                    2x-2y+2z=0, \\
                    x-z=0.
                  \end{cases}
                \]
                Persamaan pertama memberikan $z=x$.
                Substitusikan ke persamaan kedua:
                \[
                  2x-2y+2x=0\quad\Longrightarrow\quad y=2x.
                \]
                Persamaan ketiga juga terpenuhi.
                Maka,
                \[
                  v=x\begin{bmatrix}1\\2\\1\end{bmatrix},
                  \qquad
                  E_5=\Span\left\{\begin{bmatrix}1\\2\\1\end{bmatrix}\right\},
                  \qquad\dim E_5=1.
                \]

          \item Susun vektor-vektor basis ruang eigen sebagai kolom $P$.
                Pilih dua kolom dari $E_3$, lalu satu kolom dari $E_5$:
                \[
                  P=\begin{bmatrix}1&0&1\\0&1&2\\-1&0&1\end{bmatrix}.
                \]
                Ekspansi pada kolom kedua memberikan
                \[
                  \det P=\begin{vmatrix}1&1\\-1&1\end{vmatrix}=2\ne0.
                \]
                Jadi, $P$ nonsingular.
                Urutan nilai eigen mengikuti urutan kolom $P$.
                Maka, pilih
                \[
                  D=\begin{bmatrix}3&0&0\\0&3&0\\0&0&5\end{bmatrix}.
                \]
                Periksa bahwa
                \[
                  AP=\begin{bmatrix}3&0&5\\0&3&10\\-3&0&5\end{bmatrix}
                  =\begin{bmatrix}1&0&1\\0&1&2\\-1&0&1\end{bmatrix}
                  \begin{bmatrix}3&0&0\\0&3&0\\0&0&5\end{bmatrix}
                  =PD.
                \]
                Kalikan kedua ruas dengan $P^{-1}$ dari kiri.
                Diperoleh $P^{-1}AP=D$, sesuai yang diminta.
        \end{enumerate}
\end{enumerate}

\end{document}
